\section{AHSME 1991}
        \begin{enumerate}
        	\item (\textit{AHSME 1991}) If for any three distinct numbers $a$, $b$, and $c$ we define $f(a,b,c)=\frac{c+a}{c-b}$, then $f(1,-2,-3)$ is


 $\textbf{(A) } -2 \qquad \textbf{(B) } -\frac{2}{5} \qquad \textbf{(C) } -\frac{1}{4} \qquad \textbf{(D) } \frac{2}{5} \qquad \textbf {(E) } 2$


 
 \textbf{Answer: }E



\textbf{Solution}

If we plug in $1$ as $a$, $-2$ as $b$, and $-3$ as $c$ in the expression $\frac{c+a}{c-b}$, then we get $\frac{-3+1}{-3-(-2)}=\frac{-2}{-1}=2$, which is choice $\boxed{\textbf{E}}$.

 


	\item (\textit{AHSME 1991}) $|3-\pi|=$


 $\textbf{(A) }\frac{1}{7} \qquad \textbf{(B) }0.14 \qquad \textbf{(C) }3-\pi \qquad \textbf{(D) }3+\pi \qquad \textbf{(E) }\pi-3$


 
 \textbf{Answer: }E



\textbf{Solution}

Since $\pi>3$, the value of $3-\pi$ is negative. The absolute value of a negative quantity is the negative quantity multiplied by $-1$, or the negative of that quantity. Therefore $|3-\pi|=-(3-\pi)=\pi-3$, which is choice $\boxed{\textbf{E}}$

 


	\item (\textit{AHSME 1991}) $(4^{-1}-3^{-1})^{-1}=$


 $\textbf{(A) }-12 \qquad \textbf{(B) }-1 \qquad \textbf{(C) }\frac{1}{12} \qquad \textbf{(D) }1 \qquad \textbf{(E) }12$


 
 \textbf{Answer: }A



\textbf{Solution}

$\fbox{A}$ This is $\frac{1}{\frac{1}{4} - \frac{1}{3}} = \frac{1}{\frac{-1}{12}} = -12.$

 


	\item (\textit{AHSME 1991}) Which of the following triangles cannot exist?


 $\textbf{(A) }$ An acute isosceles triangle


 $\textbf{(B) }$ An isosceles right triangle


 $\textbf{(C) }$ An obtuse right triangle


 $\textbf{(D) }$ A scalene right triangle


 $\textbf{(E) }$ A scalene obtuse triangle


 
 \textbf{Answer: }C



\textbf{Solution}

$\fbox{C}$A right triangle has an angle of 90 degrees, so any obtuse angle would make the sum of those two angles over 180. This contradicts the angle sum theorem (all triangles have angles that sum to 180) since negative angles don't exist.

 


	\item (\textit{AHSME 1991}) 
                    \begin{figure}[H]
                        centering
                        \includegraphics[width=0.5\textwidth]{Problem/AHSME_I_1991_5_Problem_Figure_1.png}
                    \end{figure}
                    


 In the arrow-shaped polygon [see figure], the angles at vertices $A,C,D,E$ and $F$ are right angles, $BC=FG=5, CD=FE=20, DE=10$, and $AB=AG$. The area of the polygon is closest to


 $\textbf{(A) } 288\qquad\textbf{(B) } 291\qquad\textbf{(C) } 294\qquad\textbf{(D) } 297\qquad\textbf{(E) } 300$


 
 \textbf{Answer: }E



\textbf{Solution}

$\fbox{E}$Since they tell us that AB=AG and that angle A is a right angle, triangle ABG is a 45-45-90 triangle. Thus we can find the legs of the triangle by dividing $20$ (the hypotenuse) by $\sqrt2$. Both legs have length $10\sqrt2$, so the area of the right triangle is $100$. The rectangle is simple, just $20\times10$, so the area is 200. Adding these areas, we get $300$ as the total area.

 


	\item (\textit{AHSME 1991}) If $x\geq 0$, then $\sqrt{x\sqrt{x\sqrt{x}}}=$


 $\textbf{(A) } x\sqrt{x}\qquad \textbf{(B) } x\sqrt[4]{x}\qquad \textbf{(C) } \sqrt[8]{x}\qquad \textbf{(D) } \sqrt[8]{x^3}\qquad \textbf{(E) } \sqrt[8]{x^7}$


 
 \textbf{Answer: }E



\textbf{Solution}

Recall that square roots are one-half powers, namely $\sqrt y=y^{\frac12}$ for all $y\geq0.$

 We have
 \begin{align*} \sqrt{x\sqrt{x\sqrt{x}}} &= \sqrt{x\sqrt{x\cdot x^{\frac12}}} \\ &= \sqrt{x\sqrt{x^{\frac32}}} \\ &= \sqrt{x\cdot\left(x^{\frac32}\right)^{\frac12}} \\ &= \sqrt{x^{\frac74}} \\ &= \left(x^{\frac74}\right)^{\frac12} \\ &= x^{\frac78} \\ &= \boxed{\textbf{(E) } \sqrt[8]{x^7}}. \end{align*}
~Hapaxoromenon (Solution)

 


	\item (\textit{AHSME 1991}) If $x=\frac{a}{b}$, $a\neq b$ and $b\neq 0$, then $\frac{a+b}{a-b}=$


 $\textbf{(A) } \frac{x}{x+1} \qquad \textbf{(B) } \frac{x+1}{x-1} \qquad \textbf{(C) } 1 \qquad \textbf{(D) } x-\frac{1}{x} \qquad \textbf{(E) } x+\frac{1}{x}$


 
 \textbf{Answer: }B



\textbf{Solution}

$\frac{a+b}{a-b}= \frac{\frac{a}{b} + 1}{\frac{a}{b} - 1} = \frac{x+1}{x-1}$, so the answer is $\boxed{B}$.

 


	\item (\textit{AHSME 1991}) Liquid $X$ does not mix with water. Unless obstructed, it spreads out on the surface of water to form a circular film $0.1$cm thick. A rectangular box measuring $6$cm by $3$cm by $12$cm is filled with liquid $X$. Its contents are poured onto a large body of water. What will be the radius, in centimeters, of the resulting circular film?


 $\textbf{(A) } \frac{\sqrt{216}}{\pi} \qquad \textbf{(B) }\sqrt{\frac{216}{\pi}} \qquad \textbf{(C) } \sqrt{\frac{2160}{\pi}} \qquad \textbf{(D) } \frac{216}{\pi} \qquad \textbf{(E) } \frac{2160}{\pi}$


 
 \textbf{Answer: }C



\textbf{Solution}

$\fbox{C}$ The volume of liquid $X$ is $6 \times 3 \times 12 = 216$, so if the radius is $r$, then using the formula for the volume of a cylinder, we get $\pi r^2 \times 0.1 = 216 \implies \pi r^2 = 2160 \implies r = \sqrt{\frac{2160}{\pi}}.$

 


	\item (\textit{AHSME 1991}) From time $t=0$ to time $t=1$ a population increased by $i\%$, and from time $t=1$ to time $t=2$ the population increased by $j\%$. Therefore, from time $t=0$ to time $t=2$ the population increased by


 $\textbf{(A) } (i+j)\% \qquad \textbf{(B) } ij\% \qquad \textbf{(C) } (i+ij)\% \qquad \textbf{(D) } \left(i+j+\frac{ij}{100}\right)\% \qquad \textbf{(E) } \left(i+j+\frac{i+j}{100}\right)\%$


 
 \textbf{Answer: }D



\textbf{Solution}

The scale factors for the increases are $1+\frac{i}{100}$ and $1+\frac{j}{100}$, so the overall scale factor is $(1+\frac{i}{100})(1+\frac{j}{100}) = 1 + \frac{i}{100} + \frac{j}{100} + \frac{ij}{100^2}$. To convert this to a percentage, we subtract 1 and then multiply by 100, giving $i + j + \frac{ij}{100}$, or $\boxed{\text{D}}$.

 


	\item (\textit{AHSME 1991}) Point $P$ is $9$ units from the center of a circle of radius $15$. How many different chords of the circle contain $P$ and have integer lengths?


 $\textbf{(A) } 11\qquad \textbf{(B) } 12\qquad \textbf{(C) } 13\qquad \textbf{(D) } 14\qquad \textbf{(E) } 29$


 
 \textbf{Answer: }B



\textbf{Solution}

Let $O$ be the center of the circle, and let the chord passing through $P$ that is perpendicular to $OP$ intersect the circle at $Q$ and $R$. Then $OP = 9$ and $OQ = 15$, so by the Pythagorean Theorem, $PQ = 12$. By symmetry, $PR = 12$.

 
                    \begin{figure}[H]
                        centering
                        \includegraphics[width=0.5\textwidth]{Solution/AHSME_I_1991_10_Solution_Figure_1.png}
                    \end{figure}
                    

 


 Therefore, there are $2 + 2(29 - 25 + 1) = \boxed{12}$ chords of integer length passing through $P$.

 


	\item (\textit{AHSME 1991}) Jack and Jill run 10 km. They start at the same point, run 5 km up a hill, and return to the starting point by the same route. Jack has a 10 minute head start and runs at the rate of 15 km/hr uphill and 20 km/hr downhill. Jill runs 16 km/hr uphill and 22 km/hr downhill. How far from the top of the hill are they when they pass each other going in opposite directions (in km)?


 $\textbf{(A) } \frac{5}{4}\qquad \textbf{(B) } \frac{35}{27}\qquad \textbf{(C) } \frac{27}{20}\qquad \textbf{(D) } \frac{7}{3}\qquad \textbf{(E) } \frac{28}{49}$


 
 \textbf{Answer: }B



\textbf{Solution}

$\fbox{B}$ Consider the distance-time graph and use coordinate geometry, with time on the $x$-axis and distance on the $y$-axis. Thus Jack starts at $(0,0)$ and his initial motion, taking time in hours, is $y=15x$. This ends when $y=5$, giving the point $(\frac{1}{3}, 5)$. He now runs in the opposite direction at $20$ km/hr, so the equation is $y-5 = -20(x-\frac{1}{3})$. Now Jill starts at $(\frac{1}{6}, 0)$ (as Jack has a head start of 10 minutes = $\frac{1}{6}$ hours), so her equation is $y=16(x-\frac{1}{6}).$ Solving simultaneously with Jack's equation gives $-20x+\frac{35}{3} = 16x - \frac{8}{3} \implies \frac{43}{3} = 36x \implies x = \frac{43}{108}$, so $y = 16(\frac{25}{108}) = \frac{400}{108}$, and thus the required distance is $5 - \frac{400}{108} = \frac{140}{108} = \frac{70}{54} = \frac{35}{27}.$

 


	\item (\textit{AHSME 1991}) The measures (in degrees) of the interior angles of a convex hexagon form an arithmetic sequence of integers. Let $m$ be the measure of the largest interior angle of the hexagon. The largest possible value of $m$, in degrees, is


 $\textbf{(A) } 165\qquad \textbf{(B) } 167\qquad \textbf{(C) } 170\qquad \textbf{(D) } 175\qquad \textbf{(E) } 179$


 
 \textbf{Answer: }D



\textbf{Solution}

$\fbox{D}$ The angles must add to $180(6-2) = 720$. Now, let the arithmetic progression have first term $a$ and difference $d$, so by the sum formula, we get $\frac{6}{2}(2a+5d) = 720 \implies 2a+5d=240.$ Now, $240$ and $5d$ are both divisible by $5$ (as $d$ is an integer), so $2a$ is divisible by $5$, and thus $a$ is divisible by 5, so $m = a+5d$ is divisible by $5$. As the hexagon is convex, $m$ must be less than $180$, so as it is a multiple of $5$, it can be at most $175$, and indeed this is possible with $a = 65, d = 22.$

 


	\item (\textit{AHSME 1991}) Horses $X,Y$ and $Z$ are entered in a three-horse race in which ties are not possible. The odds against $X$ winning are $3:1$ and the odds against $Y$ winning are $2:3$, what are the odds against $Z$ winning? (By "odds against $H$ winning are $p:q$" we mean the probability of $H$ winning the race is $\frac{q}{p+q}$.)


 $\textbf{(A) } 3:20\qquad \textbf{(B) } 5:6\qquad \textbf{(C) } 8:5\qquad \textbf{(D) } 17:3\qquad \textbf{(E) } 20:3$


 
 \textbf{Answer: }D



\textbf{Solution}

Solution by e\_power\_pi\_times\_i

 
Because the odds against $X$ are $3:1$, the chance of $X$ losing is $\dfrac{3}{4}$. Since the chance of $X$ losing is the same as the chance of $Y$ and $Z$ winning, and since the odds against $Y$ are $2:3$, $Y$ wins with a probability of $\dfrac{3}{5}$. Then the chance of $Z$ winning is $\dfrac{3}{4} - \dfrac{3}{5} = \dfrac{3}{20}$. Therefore the odds against $Z$ are $\boxed{\textbf{(D) } 17:3}$.

 


	\item (\textit{AHSME 1991}) If $x$ is the cube of a positive integer and $d$ is the number of positive integers that are divisors of $x$, then $d$ could be


 $\textbf{(A) } 200\qquad \textbf{(B) } 201\qquad \textbf{(C) } 202\qquad \textbf{(D) } 203\qquad \textbf{(E) } 204$


 
 \textbf{Answer: }C



\textbf{Solution 1: Number Sense}

Solution by e\_power\_pi\_times\_i

 
Notice that if $x$ is expressed in the form $a^b$, then the number of positive divisors of $x^3$ is $3b+1$. Checking through all the answer choices, the only one that is in the form $3b+1$ is $\boxed{\textbf{(C) } 202}$.

 



\textbf{Solution 2: Answer Choices}

Solution by e\_power\_pi\_times\_i

 
Since the divisors are from $x^3$, then the answer must be something in (mod $3$). Since $200$ and $203$ are the same (mod $3$), as well as $201$ and $204$, $\boxed{\textbf{(C) } 202}$ is the only answer left.

 


	\item (\textit{AHSME 1991}) A circular table has 60 chairs around it. There are $N$ people seated at this table in such a way that the next person seated must sit next to someone. What is the smallest possible value for $N$?


 $\textbf{(A) } 15\qquad \textbf{(B) } 20\qquad \textbf{(C) } 30\qquad \textbf{(D) } 40\qquad \textbf{(E) } 58$


 
 \textbf{Answer: }B



\textbf{Solution}

$\fbox{B}$ If we fill every third chair with a person, then the condition is satisfied, giving $N=20$. Decreasing $N$ any further means there is at least one gap of $4$, so that the person can sit themselves in the middle (seat $2$ of $4$) and not be next to anyone. Hence the minimum value of $N$ is $20$.

 


	\item (\textit{AHSME 1991}) One hundred students at Century High School participated in the AHSME last year, and their mean score was 100. The number of non-seniors taking the AHSME was $50\%$ more than the number of seniors, and the mean score of the seniors was $50\%$ higher than that of the non-seniors. What was the mean score of the seniors?


 $\textbf{(A) } 100\qquad \textbf{(B) } 112.5\qquad \textbf{(C) } 120\qquad \textbf{(D) } 125\qquad \textbf{(E) } 150$


 
 \textbf{Answer: }D



\textbf{Solution}

Solution by e\_power\_pi\_times\_i

 
Let $s$ and $\dfrac{3}{2}s$ denote the numbers of seniors and non-seniors, respectively. Then $\dfrac{5}{2}s = 100$, so $s = 40$, $\dfrac{3}{2}s = 60$. Let $m$ and $\dfrac{2}{3}m$ denote the mean score of seniors and non-seniors, respectively. Then $40m + 60(\dfrac{2}{3}m) = 40m + 40m = 10000$. The answer is $\boxed{\textbf{(D) } 125}$.

 


	\item (\textit{AHSME 1991}) A positive integer $N$ is a \textit{palindrome} if the integer obtained by reversing the sequence of digits of $N$ is equal to $N$. The year 1991 is the only year in the current century with the following 2 properties:


 (a) It is a palindrome
(b) It factors as a product of a 2-digit prime palindrome and a 3-digit prime palindrome. 


 How many years in the millenium between 1000 and 2000 have properties (a) and (b)?


 $\textbf{(A) } 1\qquad \textbf{(B) } 2\qquad \textbf{(C) } 3\qquad \textbf{(D) } 4\qquad \textbf{(E) } 5$


 
 \textbf{Answer: }D



\textbf{Solution}

Solution by e\_power\_pi\_times\_i

 Notice that all four-digit palindromes are divisible by $11$, so that is our two-digit prime. Because the other factor is a three-digit number, we are looking at palindromes between $1100$ and $2000$, which also means that the last digit of the three-digit number is $1$. Checking through the three-digit numbers $101, 111, 121,\dots, 181$, we find out that there are $\boxed{\textbf{(D) } 4}$ three-digit prime numbers, which when multiplied by $11$, result in palindromes.

 


	\item (\textit{AHSME 1991}) If $S$ is the set of points $z$ in the complex plane such that $(3+4i)z$ is a real number, then $S$ is a 


 $\textbf{(A) }$ right triangle $\qquad \textbf{(B) }$ circle $\qquad \textbf{(C) }$ hyperbola $\qquad \textbf{(D) }$ line $\qquad \textbf{(E) }$ parabola


 
 \textbf{Answer: }D



\textbf{Solution}

$\fbox{D}$

 



\textbf{Solution 1}

We want $(3+4i)z$ a real number, so we want the $4i$ term to be canceled out. Then, we can make $z$ be in the form $(n-\frac{4}{3}ni)$ to make sure the imaginary terms cancel out when it's multiplied together. $(n-\frac{4}{3}ni)$ is a line, so the answer is $\textbf{(D) } \text{line}$

 



\textbf{Solution 2}

Let $z = a + bi$. Then, we have $(3 + 4i)(a + bi)$ which simplifies to $3 - 4b + (4a + 3b)i$. So whenever $4a + 3b = 0$, the value is real and thus, it produces a line. $\textbf{(D) }$~elpianista227

 


	\item (\textit{AHSME 1991}) 
                    \begin{figure}[H]
                        centering
                        \includegraphics[width=0.5\textwidth]{Problem/AHSME_I_1991_19_Problem_Figure_1.png}
                    \end{figure}
                    


 Triangle $ABC$ has a right angle at $C, AC=3$ and $BC=4$. Triangle $ABD$ has a right angle at $A$ and $AD=12$. Points $C$ and $D$ are on opposite sides of $\overline{AB}$. The line through $D$ parallel to $\overline{AC}$ meets $\overline{CB}$ extended at $E$. If

 \[\frac{DE}{DB}=\frac{m}{n},\]
where $m$ and $n$ are relatively prime positive integers, then $m+n$ is


 $\textbf{(A) } 25\qquad \textbf{(B) } 128\qquad \textbf{(C) } 153\qquad \textbf{(D) } 243\qquad \textbf{(E) } 256$


 
 \textbf{Answer: }B



\textbf{Solution 1}

Solution by e\_power\_pi\_times\_i

 
Let $F$ be the point such that $DF$ and $CF$ are parallel to $CE$ and $DE$, respectively, and let $DE = x$ and $BE^2 = 169-x^2$. Then, $[FDEC] = x(4+\sqrt{169-x^2}) = [ABC] + [BED] + [ABD] + [AFD] = 6 + \dfrac{x\sqrt{169-x^2}}{2} + 30 + \dfrac{(x-3)(4+\sqrt{169-x^2})}{2}$. So, $4x+x\sqrt{169-x^2} = 60 + x\sqrt{169-x^2} - 3\sqrt{169-x^2}$. Simplifying $3\sqrt{169-x^2} = 60 - 4x$, and $1521 - 9x^2 = 16x^2 - 480x + 3600$. Therefore $25x^2 - 480x + 2079 = 0$, and $x = \dfrac{48\pm15}{5}$. Checking, $x = \dfrac{63}{5}$ is the answer, so $\dfrac{DE}{DB} = \dfrac{\dfrac{63}{5}}{13} = \dfrac{63}{65}$. The answer is $\boxed{\textbf{(B) } 128}$.

 



\textbf{Solution 3 (Trig)}

We have $\angle ABC = \arcsin\left(\frac{3}{5}\right)$ and $\angle DBA=\arcsin\left(\frac{12}{13}\right).$ Now we are trying to find $\sin(\angle DBE)=\sin\left((180^{\circ}-\angle DBC)\right)=\sin\left(180^{\circ}-\arcsin\left(\frac{3}{5}\right)-\arcsin\left(\frac{12}{13}\right)\right)=\sin(\arcsin\left(\frac{3}{5}\right)+\arcsin\left(\frac{12}{13}\right)).$ Now we use the $\sin$ angle sum identity, which states $\sin(a+b)=\sin(a)\cos(b)+\cos(a)\sin(b).$ Using this identity yields $\sin\left(\arcsin\left(\frac{3}{5}\right)\right)\cos\left(\arcsin\left(\frac{12}{13}\right)\right)+\cos\left(\arcsin\left(\frac{3}{5}\right)\right)\sin\left(\arcsin\left(\frac{12}{13}\right)\right).$

 


	\item (\textit{AHSME 1991}) The sum of all real $x$ such that $(2^x-4)^3+(4^x-2)^3=(4^x+2^x-6)^3$ is


 $\textbf{(A) } \frac32 \qquad \textbf{(B) } 2 \qquad \textbf{(C) } \frac52 \qquad \textbf{(D) } 3 \qquad \textbf{(E) } \frac72$


 
 \textbf{Answer: }E



\textbf{Solution}

Note that $(2^x-4)+(4^x-2)=4^x+2^x-6,$ so we let $a=2^x-4$ and $b=4^x-2.$ The original equation becomes 
 \[a^3+b^3=(a+b)^3.\]We expand the right side, then rearrange:
 \begin{align*} a^3+b^3 &= a^3+3a^2b+3ab^2+b^3 \\ 0 &= 3a^2b+3ab^2 \\ 0 &= 3ab(a+b). \end{align*}


 <ul style="list-style-type:square;"> \item If $a=0,$ then $2^x-4=0,$ from which $x=2.$ 

 \item If $b=0,$ then $4^x-2=0,$ from which $x=\frac12.$ 

 <li>If $a+b=0,$ then $4^x+2^x-6=0.$ As $4^x=(2^x)^2,$ we rewrite this equation, then factor:
 \begin{align*} (2^x)^2+2^x-6&=0 \\ (2^x-2)(2^x+3)&=0. \end{align*}
If $2^x-2=0,$ then $x=1.$ If $2^x+3=0,$ then there are no real solutions for $x,$ as $2^x+3>3$ holds for all real numbers $x.$

 </li>

 </ul>Together, the answer is $2+\frac12+1=\boxed{\textbf{(E) } \frac72}.$

 


	\item (\textit{AHSME 1991}) For all real numbers $x$ except $x=0$ and $x=1$ the function $f(x)$ is defined by $f(x/(x-1))=1/x$. Suppose $0\leq t\leq \pi/2$. What is the value of $f(\sec^2t)$?


 $\textbf{(A) } \sin^2\theta\qquad \textbf{(B) } \cos^2\theta\qquad \textbf{(C) } \tan^2\theta\qquad \textbf{(D) } \cot^2\theta\qquad \textbf{(E) } \csc^2\theta$


 
 \textbf{Answer: }A



\textbf{Solution}

Let $y=\frac{x}{x-1} \Rightarrow xy-y=x \Rightarrow x=\frac{y}{y-1}$

 $f(y)=\frac{1}{x}=\frac{y-1}{y}=1-\frac{1}{y}$

 $f(\sec^2t)=\sin^2t$

 $\fbox{A}$

 


	\item (\textit{AHSME 1991}) 
                    \begin{figure}[H]
                        centering
                        \includegraphics[width=0.5\textwidth]{Problem/AHSME_I_1991_22_Problem_Figure_1.png}
                    \end{figure}
                    


 Two circles are externally tangent. Lines $\overline{PAB}$ and $\overline{PA'B'}$ are common tangents with $A$ and $A'$ on the smaller circle $B$ and $B'$ on the larger circle. If $PA=AB=4$, then the area of the smaller circle is


 $\textbf{(A) } 1.44\pi\qquad \textbf{(B) } 2\pi\qquad \textbf{(C) } 2.56\pi\qquad \textbf{(D) } \sqrt{8}\pi\qquad \textbf{(E) } 4\pi$


 
 \textbf{Answer: }B



\textbf{Solution}

Using the tangent-tangent theorem, $PA=AB=PA'=A'B'=4$. We can then drop perpendiculars from the centers of the circles to the points of tangency and use similar triangles. We let the radius of the larger circle be $x$ and the radius of the smaller circle be $y$, so we can see that, using similar triangles, $x=2y$. In addition, the total hypotenuse of the larger right triangles equals $2(x+y)$ since half of it is $x+y$, so $y^2+4^2=(3y)^2$. If we simplify, we get $y^2+16=9y^2$, so $8y^2=16$, so $y=\sqrt2$. This means that the smaller circle has area $2\pi$, which is answer choice $\fbox{B}$.

 


	\item (\textit{AHSME 1991}) 
                    \begin{figure}[H]
                        centering
                        \includegraphics[width=0.5\textwidth]{Problem/AHSME_I_1991_23_Problem_Figure_1.png}
                    \end{figure}
                    


 If $ABCD$ is a $2\times2$ square, $E$ is the midpoint of $\overline{AB}$,$F$ is the midpoint of $\overline{BC}$,$\overline{AF}$ and $\overline{DE}$ intersect at $I$, and $\overline{BD}$ and $\overline{AF}$ intersect at $H$, then the area of quadrilateral $BEIH$ is


 $\textbf{(A) } \frac{1}{3}\qquad \textbf{(B) } \frac{2}{5}\qquad \textbf{(C) } \frac{7}{15}\qquad \textbf{(D) } \frac{8}{15}\qquad \textbf{(E) } \frac{3}{5}$


 
 \textbf{Answer: }C



\textbf{Solution 1: Coordinate Geometry}

Solution by you\_r\_gassy

 
First, we find out the coordinates of the vertices of quadrilateral $BEIH$, then use the Shoelace Theorem to solve for the area. Denote $B$ as $(0,0)$. Then $E (0,1)$. Since I is the intersection between lines $DE$ and $AF$, and since the equations of those lines are $y = \dfrac{1}{2}x + 1$ and $y = -2x + 2$, $I (\dfrac{2}{5}, \dfrac{6}{5})$. Using the same method, the equation of line $BD$ is $y = x$, so $H (\dfrac{2}{3}, \dfrac{2}{3})$. Using the Shoelace Theorem, the area of $BEIH$ is $\dfrac{1}{2}\cdot\dfrac{14}{15} = \boxed{\textbf{(C) } \dfrac{7}{15}}$.

 


	\item (\textit{AHSME 1991}) The graph, $G$ of $y=\log_{10}x$ is rotated $90^{\circ}$ counter-clockwise about the origin to obtain a new graph $G'$. Which of the following is an equation for $G'$?


 $\textbf{(A) } y=\log_{10}\left(\frac{x+90}{9}\right) \qquad \textbf{(B) } y=\log_{x}10 \qquad \textbf{(C) } y=\frac{1}{x+1} \qquad \textbf{(D) } y=10^{-x} \qquad \textbf{(E) } y=10^x$


 
 \textbf{Answer: }D



\textbf{Solution}

$\fbox{D}$ Rotating a point $(x,y)$ $90^{\circ}$ anticlockwise about the origin maps it to $(-y,x).$ (You can prove this geometrically, or using matrices if you aren't convinced). Thus $(x, \log_{10}x)$ maps to $(-\log_{10}x, x)$, so new $y =$ old $x = 10^{\log_{10}x} = 10^{-(-\log_{10}x)} = 10^{-\text{new} x}$, so the new equation is $y=10^{-x}.$

 


	\item (\textit{AHSME 1991}) If $T_n=1+2+3+\cdots +n$ and

 \[P_n=\frac{T_2}{T_2-1}\cdot\frac{T_3}{T_3-1}\cdot\frac{T_4}{T_4-1}\cdot\cdots\cdot\frac{T_n}{T_n-1}\]
for $n=2,3,4,\cdots,$ then $P_{1991}$ is closest to which of the following numbers?


 $\textbf{(A) } 2.0\qquad \textbf{(B) } 2.3\qquad \textbf{(C) } 2.6\qquad \textbf{(D) } 2.9\qquad \textbf{(E) } 3.2$


 
 \textbf{Answer: }D



\textbf{Solution}

$\fbox{D}$ It is well known that $T_n = \frac{n(n+1)}{2}$, so we can put this in and simplify to get that the general term of the product is $\frac{n(n+1)}{(n+2)(n-1)}.$ Now by writing out the terms, we see that almost everything cancels (telescopes), giving $\frac{3 \times 1991}{1993}$ which is almost $3$, so closest to $2.9.$

 


	\item (\textit{AHSME 1991}) An $n$-digit positive integer is cute if its $n$ digits are an arrangement of the set $\{1,2,...,n\}$ and its first 
$k$ digits form an integer that is divisible by $k$  , for  $k  = 1,2,...,n$. For example, $321$ is a cute $3$-digit integer because $1$ divides $3$, $2$ divides $32$, and $3$ divides $321$. How many cute $6$-digit integers are there?


 $\textbf{(A) } 0\qquad \textbf{(B) } 1\qquad \textbf{(C) } 2\qquad \textbf{(D) } 3\qquad \textbf{(E) } 4$


 
 \textbf{Answer: }C



\textbf{Solution 1}

Suppose $ABCDEF$ is a $6$-digit cute integer, where $A$, $B$, $C$, $D$, $E$, and $F$ are different digits between $1$ and $6$ inclusive.

 Since $ABCDE$ is divisible by $5$, $E = 5$.

 Since $AB$ is divisible by $2$, $ABCD$ is divisible by $4$, and $ABCDEF$ is divisible by $6$, all three of these numbers are even.  So $\{B, D, F\} = \{2, 4, 6\}$, and thus $\{A, C\} = \{1, 3\}$.

 Since $ABC$ is divisible by $3$, $A + B + C = 4 + B$ is also divisible by $3$.  Thus $B = 2$ and $\{D, F\} = \{4, 6\}$.

 Since $ABCD$ is divisible by $4$, and $C \in \{1, 3\}$, $D = 6$ and thus $F = 4$.

 Thus $ABCDEF = A2C654$, with $\{A, C\} = \{1, 3\}$.  Both $123654$ and $321654$ are cute, and these must be the only $\boxed{\textbf{(C)}\ 2}$ cute $6$-digit numbers. 

 -j314andrews 

 



\textbf{Solution 2}

$\fbox{C}$ Let the number be $abcdef$. We know $1$ will always divide $a$. $5$ must divide $abcde$, so $e$ must be $5$ or $0$, but we can only use the digits $1$ to $6$, so $e = 5$. $4$ must divide $abcd$, so it must divide $cd$ (the test for divisibility by 4 is that the last two digits form a number divisible by 4), and so $cd$ must be $12$, $16$, $24$, $32$, $36$, (not $52$ or $56$ as the $5$ is already used), or $64$. We know $2$ divides $ab$ so $b$ is even, $4$ divides $abcd$ so $d$ is even, and $6$ divides $abcdef$ so $f$ is even, and thus $b$, $d$, and $f$ must be $2$, $4$, and $6$ in some order, so $c$ must be odd, so $cd$ must be $12$, $16$, $32$, or $36$. Now $6$ divides $abcdef$ implies $3$ divides $abcdef$, and we know $3$ divides $abc$, so $3$ must also divide $def$, so we need a multiple of $3$ starting with $2$ or $6$, using the remaining digits. If $cd$ is $12$, then $def$ must be $243$, $234$, $246$, or $264$, but none of these work as we need $e = 5$. If $cd$ is $16$, $def$ must be $624$, $642$, $645$, or $654$, so $654$ works (it has $e=5$). If $cd$ is $32$, similarly nothing works, and if $cd$ is $36$, $651$ or $654$ works. So now we have, as possibilities, $ab1654$, $ab3651$, and $ab3654$. But clearly $ab3651$ can't work, as we need $abcdef$ to be divisible by $6$, so it must be even, so we eliminate this possibility. Now with $ab1654$, we need $b$ to be even, so it must be $2$, giving $a21654$, and then $321654$ works as $3$ does divide $321$. With $ab3654$, we get $a23654$ as $b$ is even, and then $123654$ works as $3$ divides $123$. Hence the number of such numbers is $2$: $321654$ and $123654$.

 


	\item (\textit{AHSME 1991}) If 
 \[x+\sqrt{x^2-1}+\frac{1}{x-\sqrt{x^2-1}}=20,\] then 
 \[x^2+\sqrt{x^4-1}+\frac{1}{x^2+\sqrt{x^4-1}}=\]


 $\textbf{(A) } 5.05 \qquad \textbf{(B) } 20 \qquad \textbf{(C) } 51.005 \qquad \textbf{(D) } 61.25 \qquad \textbf{(E) } 400$


 
 \textbf{Answer: }C



\textbf{Solution}

We rationalize the denominator in the given equation, then solve for $x:$
 \begin{align*} x+\sqrt{x^2-1}+\frac{x+\sqrt{x^2-1}}{\left(x-\sqrt{x^2-1}\right)\left(x+\sqrt{x^2-1}\right)} &= 20 \\ x+\sqrt{x^2-1}+x+\sqrt{x^2-1} &= 20 \\ x+\sqrt{x^2-1} &= 10 \\ \sqrt{x^2-1} &= 10-x \\ x^2-1 &= 100-20x+x^2 \\ 20x &= 101 \\ x &= 5.05. \end{align*}
We rationalize the denominator in the requested expression, then simplify the result:
 \begin{align*} x^2+\sqrt{x^4-1}+\frac{1}{x^2+\sqrt{x^4-1}} &= x^2+\sqrt{x^4-1}+\frac{x^2-\sqrt{x^4-1}}{\left(x^2+\sqrt{x^4-1}\right)\left(x^2-\sqrt{x^4-1}\right)} \\ &= x^2+\sqrt{x^4-1}+x^2-\sqrt{x^4-1} \\ &= 2x^2 \\ &= \boxed{\textbf{(C) } 51.005}. \end{align*}


 


	\item (\textit{AHSME 1991}) Initially an urn contains 100 white and 100 black marbles. Repeatedly 3 marbles are removed (at random) from the urn and replaced with some marbles from a pile outside the urn as follows: 3 blacks are replaced with 1 black, or 2 blacks and 1 white are replaced with a white and a black, or 1 black and 2 whites are replaced with 2 whites, or 3 whites are replaced with a black and a white. Which of the following could be the contents of the urn after repeated applications of this procedure?


 $\textbf{(A) }$ 2 black $\qquad \textbf{(B) }$ 2 white $\qquad \textbf{(C) }$ 1 black $\qquad \textbf{(D) }$ 1 black and 1 white $\qquad \textbf{(E) }$ 1 white


 
 \textbf{Answer: }B



\textbf{Solution}

The four operations result in the following net changes:

 
                \begin{itemize}
                    \item Replacing $3$ black marbles with $1$ black marble results in a loss of $2$ black marbles.\item Replacing $2$ black marbles and $1$ white marble with $1$ black marble and $1$ white marble results in a loss of $1$ black marble.\item Replacing $1$ black marble and $2$ white marbles with $2$ white marbles results in a loss of $1$ black marble.\item Replacing $3$ white marbles with $1$ white marble with $2$ black marbles results in a loss of $2$ white marbles and a gain of $1$ black marble.
                \end{itemize}


                Notice that only the fourth operation changes the number of white marbles.  Since this operation reduces the number of white marbles by $2$ each time, the number of white marbles must always be even.  This eliminates choices (D) and (E).

 Also, notice that any operation which removes at least one white marble adds at least one white marble back to the urn.  Therefore, it is impossible to remove all the white marbles.  This eliminates choices (A) and (C).

 Therefore, only $\boxed{\textbf{(B)}\text{ 2 white}}$ marbles could be the contents of the urn at some point. If the fourth operation is performed $49$ times and then the third operation is performed $149$ times, the urn will contain $2$ white marbles.

 


 


	\item (\textit{AHSME 1991}) Equilateral triangle $ABC$ has $P$ on $AB$ and $Q$ on $AC$. The triangle is folded along $PQ$ so that vertex $A$ now rests at $A'$ on side $BC$. If $BA'=1$ and $A'C=2$ then the length of the crease $PQ$ is


 $\textbf{(A) } \frac{8}{5} \qquad \textbf{(B) } \frac{7}{20}\sqrt{21} \qquad \textbf{(C) } \frac{1+\sqrt{5}}{2} \qquad \textbf{(D) } \frac{13}{8} \qquad \textbf{(E) } \sqrt{3}$


 
 \textbf{Answer: }B



\textbf{Solution}

$ABC$ has side length $3$. Let $AP=A'P=x$ and $AQ=A'Q=y$. Thus, $BP=3-x$ and $CQ=3-y$. Applying Law of Cosines on triangles $BPA'$ and $CQA'$ using the $60^{\circ}$ angles gives $x=\frac{7}{5}$ and $y=\frac{7}{4}$. Applying Law of Cosines once again on triangle $APQ$ using the $60^{\circ}$ angle gives 
 \[PQ^2=\frac{(21)(49)}{400}\] so 
 \[PQ=\frac{7}{20}\sqrt{21}\] The correct answer is $\fbox{(B)}$.

 


	\item (\textit{AHSME 1991}) For any set $S$, let $|S|$ denote the number of elements in $S$, and let $n(S)$ be the number of subsets of $S$, including the empty set and the set $S$ itself. If $A$, $B$, and $C$ are sets for which $n(A)+n(B)+n(C)=n(A\cup B\cup C)$ and $|A|=|B|=100$, then what is the minimum possible value of $|A\cap B\cap C|$?


 $\textbf{(A) } 96 \qquad \textbf{(B) } 97 \qquad \textbf{(C) } 98 \qquad \textbf{(D) } 99 \qquad \textbf{(E) } 100$


 
 \textbf{Answer: }B



\textbf{Solution 1}

For any set $S$, $n(S) = 2^{|S|}$.  So $n(A)=n(B)=2^{100}$, and thus $2^{101} + 2^{|C|} = 2^{|A \cup B \cup C|}$. 

 Since $2^{101} + 2^{|C|}$ must be a power of $2$, $|C| = 101$ and thus $|A \cup B \cup C| = 102$.

 Of the $102$ elements of $A \cup B \cup C$, $2$ are not in $A$, $2$ are not in $B$, and $1$ is not in $C$.  Therefore, at most $2 + 2 + 1 = 5$ elements of $A \cup B \cup C$ are not in $A \cap B \cap C$, so $|A \cap B \cap C| \geq 102 - 5 = 97$.  One case in which $|A \cap B \cap C| = \boxed{\textbf{(B)}\ 97}$ is when $A$ is the set of all integers from $3$ to $102$, $B$ is the set of all integers from $1$ to $102$ except $3$ and $4$, and $C$ is the set of all integers from $1$ to $102$ except $5$.

 



\textbf{Solution 2 (PIE)}

As $|A|=|B|=100$, $n(A)=n(B)=2^{100}.$

 Because we're given that $n(A)+n(B)+n(C)=n(A \cup B \cup C)$, we know that$2^{|A|}+2^{|B|}+2^{|C|}=2^{|A \cup B \cup C|}$, so we can write 

 
 \[2^{100}+2^{100}+2^{|C|}=2^{|A \cup B \cup C|}\]
 \[\Rightarrow 2^{101}+2^{|C|}=2^{|A \cup B \cup C|}\]

 Because $|C|$ and $|A \cup B \cup C|$ are integers, $|C|=101$ and $|A \cup B \cup C| = 102.$

 
By the \href{https://artofproblemsolving.com/wiki/index.php?title=Principle\_of\_Inclusion-Exclusion}{Principle of Inclusion-Exclusion}, $|A \cup B| = |A| + |B| - |A \cap B| = 200 - |A \cap B|.$

 Since $|A|=|B| \le |A \cup B| \le |A \cup B \cup C|$, we know $100 \le |A \cup B| \le 102$, so $98 \le |A \cap B| \le 100.$

 
By the \href{https://artofproblemsolving.com/wiki/index.php?title=Principle\_of\_Inclusion-Exclusion}{Principle of Inclusion-Exclusion}, $|A \cup C| = |A| + |C| - |A \cap C| = 201 - |A \cap C|.$

 Since $|C| \le |A \cup C| \le |A \cup B \cup C|$, we know $101 \le |A \cup C| \le 102$, so $99 \le |A \cap C| \le 100.$

 
By the \href{https://artofproblemsolving.com/wiki/index.php?title=Principle\_of\_Inclusion-Exclusion}{Principle of Inclusion-Exclusion}, $|B \cup C| = |B| + |C| - |B \cap C| = 201 - |B \cap C|.$

 Since $|C| \le |B \cup C| \le |A \cup B \cup C|$, we know $101 \le |B \cup C| \le 102$,so $99 \le |B \cap C| \le 100.$

 
By the \href{https://artofproblemsolving.com/wiki/index.php?title=Principle\_of\_Inclusion-Exclusion}{Principle of Inclusion-Exclusion}, 

 \begin{align}|A \cap B \cap C| &= |A \cup B \cup C| - |A| - |B| - |C| + |A \cap B| + |A \cap C| + |B \cap C| \\&= 102 - 100 - 100 - 101 + |A \cap B| + |A \cap C| + |B \cap C| \\&= |A \cap B| + |A \cap C| + |B \cap C| - 199\end{align}

 
Therefore,
 \[98 + 99 + 99 - 199 \le |A \cap B \cap C| \le 100+100+100-199\]

 
 \[\Rightarrow \boxed{\textbf{97}} \le |A \cap B \cap C| \le 101.\]

 



\textbf{Solution 3}

Represent the elements of $A, B, C$ as an ordered $102$-tuple of $0$'s and $1$'s. $A$ and $B$ contain exactly $100$ $1$'s, while $C$ contains $101$ $1$'s. We want to minimize the number of $3$'s after summing the numbers in the respective positions of these $102$-tuples. In the most optimal situation, all positions of the resultant tuple contains at least a $2$; this leaves $100+100+101-2\times102=97$ positions with a $3$. Thus, the minimum value is $\boxed{B: 97}$, with construction given in the above solutions.

 -cretinouscretin

 


\end{enumerate}